\subsection{Other Properties Preserved by the Morita Correspondence}

Let A and B be Morita equivalent k-algebras and P an invertible \((A,B)_{k}\) -bimodule.

PROPOSITION 11. --- Let

\[
\mathrm{V} ^ {\prime} \xrightarrow {f} \mathrm{V} \xrightarrow {g} \mathrm{V} ^ {\prime \prime}\tag{ε}
\]

be a diagram of B-modules and B-linear mappings, and let

\[
(\mathrm{P} \otimes \mathcal {E})
\]

\[
\mathrm{P} \otimes_ {\mathrm{B}} \mathrm{V} ^ {\prime} \xrightarrow {1 _ {\mathrm{P}} \otimes f} \mathrm{P} \otimes_ {\mathrm{B}} \mathrm{V} \xrightarrow {1 _ {\mathrm{P}} \otimes g} \mathrm{P} \otimes_ {\mathrm{B}} \mathrm{V} ^ {\prime \prime}
\]

be the corresponding diagram of A-modules. Then \((\mathcal{E})\) is an exact sequence if and only if \((\mathrm{P} \otimes \mathcal{E})\) is an exact sequence.

Suppose that the sequence \((\mathcal{E})\) is exact. Since the right B-module P is projective, the sequence \((P \otimes \mathcal{E})\) is exact (II, §3, No.~6, p.~251, Proposition 5 and II, §3, No.~7, p.~257, Corollary 6).

Conversely, suppose that the sequence \((\mathrm{P} \otimes \mathcal{E})\) is exact. Let Q be a \((\mathrm{B}, \mathrm{A})_{k}\) -bimodule, inverse to P, and \(\theta : Q \otimes_{A} P \to {}_{s} B_{d}\) an isomorphism. Consider the commutative diagram

\[
\begin{array}{c} \mathrm{Q} \otimes_ {\mathrm{A}} \mathrm{P} \otimes_ {\mathrm{B}} \mathrm{V} ^ {\prime} \xrightarrow {1 _ {\mathrm{Q}} \otimes 1 _ {\mathrm{P}} \otimes f} \mathrm{Q} \otimes_ {\mathrm{A}} \mathrm{P} \otimes_ {\mathrm{B}} \mathrm{V} \xrightarrow {1 _ {\mathrm{Q}} \otimes 1 _ {\mathrm{P}} \otimes g} \mathrm{Q} \otimes_ {\mathrm{A}} \mathrm{P} \otimes_ {\mathrm{B}} \mathrm{V} ^ {\prime \prime} \\ \theta \otimes 1 _ {\mathrm{V} ^ {\prime}} \Biggl \downarrow \\ \mathrm{V} ^ {\prime} \xrightarrow {f} \mathrm{V} \xrightarrow {g} \mathrm{V} ^ {\prime \prime}. \end{array}
\]

Since Q is a projective A-module and the sequence \((\mathrm{P} \otimes \mathcal{E})\) is exact, the first line of this diagram is an exact sequence. Since the vertical arrows are isomorphisms, the second line is also exact.

COROLLARY. --- Let \(f: V \rightarrow W\) be a B-linear mapping. Then f is injective (resp. surjective) if and only if \(1_{P} \otimes f\) is.

PROPOSITION 12. --- Let V be a left B-module. The B-module V is projective (resp. generating, faithful, *injective, finitely presented*) if and only if the A-module P ⊗\_B V is.

\begin{enumerate}
\def\labelenumi{\alph{enumi})}
\item
  Suppose that V is projective. There exists a set I such that V is isomorphic to a direct factor submodule of \(\mathrm{B}_{s}^{(\mathrm{I})}\) . The A-module \(P \otimes_{B} V\) is then isomorphic to a direct factor submodule of \(\mathrm{P}^{(\mathrm{I})}\) ; since P is a projective A-module, the same holds for \(P \otimes_{B} V\) .
\item
  Suppose that the B-module V is generating. Let M be an A-module. There exists a B-module W such that M is isomorphic to \(P \otimes_{B} W\) . By Theorem 1 of VIII, p.~80, there exist a set I and a surjection \(\varphi : V^{(I)} \to W\) . By the corollary, the mapping \(1_{P} \otimes \varphi\) from \(P \otimes (V^{(I)})\) to \(P \otimes_{A} W\) is surjective, which gives a surjection \((P \otimes V)^{(I)} \to M\) . By Theorem 1 of VIII, p.~80, \(P \otimes V\) is a generating A-module.
\item
  The B-module V is faithful if and only if its annihilator is reduced to 0. Assertion c) therefore follows from Example 3 of VIII, p.~105.
\end{enumerate}

*d) Suppose that V is injective. By the remark of VIII, p.~106, the A-module \(\mathrm{P} \otimes_{\mathrm{B}} \mathrm{V}\) is isomorphic to \(\mathrm{Hom}_{\mathrm{B}}(\mathrm{Q}, \mathrm{V})\) , where Q is a \((\mathrm{B}, \mathrm{A})_k\) -bimodule inverse to A. Since the A-module Q is projective, hence flat (X, §1, n° 3, p.~9, example 1), the A-module \(\mathrm{Hom}_{\mathrm{B}}(\mathrm{Q}, \mathrm{V})\) is injective by X, §1, n° 8, p.~18, proposition 11.

\begin{enumerate}
\def\labelenumi{\alph{enumi})}
\setcounter{enumi}{4}
\item
  Suppose that V admits a finite presentation \(L_{1} \to L_{0} \to V \to 0\) (X, §1, n° 4, p.~10). By taking the tensor product with P, we deduce an exact sequence of A-modules \(N_{1}^{\prime} \xrightarrow{u} N_{0}^{\prime} \to P \otimes_{B} V \to 0\) , where \(N_{1}^{\prime}\) and \(N_{0}^{\prime}\) are projective and finitely generated (Proposition 11 and a)). Let \(N_{0}^{\prime\prime}\) be a finitely generated A-module such that the module \(N_{0} = N_{0}^{\prime} \oplus N_{0}^{\prime\prime}\) is free and finitely generated, and let \(u^{\prime}: N_{1}^{\prime} \oplus N_{0}^{\prime\prime} \to N_{0}\) be the homomorphism \((u, 1_{N_{0}^{\prime\prime}})\) ; then \(P \otimes_{B} V\) can be identified with the cokernel of \(u^{\prime}\) . Let \(N_{1}\) be a finitely generated free A-module and \(p: N_{1} \to N_{1}^{\prime} \oplus N_{0}^{\prime\prime}\) a surjective homomorphism; the sequence \(N_{1} \xrightarrow{u^{\prime} \circ p} N_{0} \to P \otimes_{B} V \to 0\) is a finite presentation of the A-module \(P \otimes_{B} V_*\)
\item
  Suppose that the A-module \(P \otimes_{B} V\) is projective (resp. generating, faithful, *injective, finitely presented*). By applying the above (interchanging the roles of A and B and of P and Q), we see that the B-module \(Q \otimes_{A} P \otimes_{B} V\) also has this property. The same is therefore true for the B-module V, which is isomorphic to it.
\end{enumerate}

COROLLARY. --- The ring A is left Artinian (resp. left Noetherian) if and only if the ring B is.

Because of the isomorphism between the ordered set of left ideals of B and the set of A-submodules of P, the ring B is left Artinian (resp. left Noetherian) if and only if the A-module P is Artinian (resp. Noetherian). However, by Theorem 1 of VIII, p.~101, the A-module P is generating and finitely generated; in particular, \(A_{s}\) is isomorphic to a direct factor of \(P^{n}\) for an integer \(n \geqslant 1\) . Consequently, P is Artinian (resp. Noetherian) if and only if A is left Artinian (resp. left Noetherian).

